\section{Problem and simulation model}
\subsection{A receiver-side mission objective}
Let $p_i(t)\in\R^2$ be the position of vessel $i=1,\ldots,N$, with a base at
the origin. Monitoring points $x_m$ sample three concentric annuli at
$0.9R,R,1.1R$, with normalized priority weights $w_m(t)$. This is quadrature of
an \emph{annular monitoring distribution}; it is not an estimate of Cartesian
area coverage over an entire disk.

Let $\tau_m^b(t)$ be the source timestamp of the newest report for point $m$
available at the base. Define
\begin{align}
 A_m(t)&=t-\tau_m^b(t),\\
 C(t)&=\sum_m w_m(t)\1\{A_m(t)\leq D\},\label{eq:C}\\
 \bar C&=\frac{1}{T-t_w}\int_{t_w}^{T}C(t)\,dt .
\end{align}
Never-received reports have infinite age. Actual implementation averages the
discrete samples with $t\geq t_w$, including the endpoint at warm-up. We also
report AoI clipped at $10D=300$ s, its saturation fraction, and fleet energy
$E=\sum_i\int_0^T P_i(t)\,dt/3600$ in Wh. Deployment energy is included.
The primary metric concerns delivered freshness, not guaranteed target tracking,
classification correctness, or collision safety.

\subsection{Sensing and environmental dependence}
A synthetic severity $s(t)\in[0,1]$ changes nominal sensing ranges and visibility.
It is not a calibrated sea-state number. The nominal radar half-range is
$r_r(s)=1650\exp(-1.4s)$ m and the camera half-range is
$r_c(s)=650\exp(-0.8s)$ m. Radar range additionally depends on bearing relative
to wind and an unknown vessel-specific multiplier $\theta_i$.
For logistic roll-off $\sigma$, radar conditional detection is
\begin{equation}
 p^r_{im}=0.94\,\sigma\!\left(\frac{r_{im}-d_{im}}{0.2r_{im}}\right),
\end{equation}
with an analogous camera curve of peak 0.86 and relative width 0.22.
Conditional sensor fusion is $p_{im}=1-(1-p^r_{im})(1-p^c_{im})$.
A single visibility variable $V_i(t)$ gates \emph{both} sensors on a vessel;
their obscuration failures are not treated as independent.

Visibility follows a correlated Gaussian field with a shared spatial component
and vessel-specific response, transformed to a Bernoulli variable with
availability
\begin{equation}
 a_i(s)=\sigma\!\left(\operatorname{logit}
 \frac{0.94}{1+1.1s^2}+b_i\right).
\end{equation}
The simulator uses a random Fourier spatial field and a 15 s default temporal
memory. The planner approximates its spatial covariance by a radial-basis kernel.
Detector noise remains independent conditional on the shared visibility paths.
Thus even with correct nominal curves, the planning approximation is not an
exact physical posterior.

\subsection{Radio paths, cache causality and estimation}
For two nodes separated by $d_{ij}$, the nominal reciprocal link probability is
\begin{equation}
 q_{ij}=\frac{0.995}{1+(d_{ij}/r_{50}(s))^6},\qquad
 r_{50}(s)=r_{50}(0)e^{-1.05s}.\label{eq:q}
\end{equation}
An episode-specific unknown radio multiplier modifies the actual process.
A common Markov fade multiplies all off-diagonal link probabilities by either
1 or 0.25. Realized links are conditionally independent given the fade, but
their marginal failures are correlated. Full blackout overrides the channel.

Each five-second communication round forwards the newest cached record across
one realized edge. Reports, health/pose telemetry and formation commands all
use this rule. A record cannot traverse two hops during one round.
The first hop can occur in the report-generation round; further hops require
subsequent rounds. The model omits finite packet capacity, interference-aware
scheduling, asymmetric links and MAC contention.

Each vessel estimates radar-range multiplier and visibility offset on a
two-dimensional Bayesian grid using fresh known reference-vessel beacons.
A missing beacon is not classified as a radar miss. Joint estimation is
challenging because poor detection can indicate either reduced range or
reduced availability. Process diffusion and a change hazard prevent permanent
overconfidence. The coordinator uses only received cached beliefs and poses.
The new controllers use posterior means in the detection predictor; they do
not integrate the full health posterior or impose an uncertainty chance
constraint. Stale-state suppression is a heuristic.

\subsection{Temporal delivery prediction}
For a prospective geometry, sample $S$ joint paths of visibility and radio graphs
over $L=\lfloor D/\Delta t\rfloor+1$ rounds. Let $R_{ik}$ indicate whether a
report generated by vessel $i$ at round $k$ can reach the base by the common
deadline, computed by backward reachability with waiting in caches.
Conditional on one such path $\omega$, the detector-noise-integrated score is
\begin{equation}
 F(\omega,p)=\sum_m w_m
 \left[1-\prod_{k=1}^{L}\prod_{i=1}^{N}
 (1-p_{im}V_{ik}R_{ik})\right].\label{eq:temporal}
\end{equation}
The predictor $\widehat F=S^{-1}\sum_\omega F(\omega,p)$ uses 32 scrambled Sobol
paths by default, independent of simulation randomness. It retains shared
visibility, shared channel fade, shared routes and path ordering. It does not
multiply independent marginal path successes. The equality in
(\ref{eq:temporal}) is conditional on the specified path model and independent
detector noise; it is not a claim that the model matches a particular radar.
Geometry is frozen inside each 30 s delivery-prediction window. Evaluating this
score at three future positions approximates movement over the longer horizon;
it does not exactly integrate changing geometry inside every report window.
Candidates share the planning samples to reduce comparison noise. Selecting
the best candidate on that same finite sample bank can nevertheless produce
selection optimism. There is no independent rollout validation bank in this
release, and the numerical acceptance margin is not a confidence bound.

\section{Economic formation adaptation}
\subsection{Finite movement forecast}
The economic controller considers small waypoint changes rather than only
switching among large predefined formations. For initial waypoint error
$d_i(0)=\|g_i-p_i\|$, it forecasts
\begin{equation}
 \dot d_i=-\min(kd_i,v_{\max}),\qquad k=0.025~\mathrm{s}^{-1}.
\end{equation}
Set $t_i^c=[d_i(0)-v_{\max}/k]_+/v_{\max}$ and
$d_i^e=\min(d_i(0),v_{\max}/k)$. Then
\begin{equation}
 d_i(t)=d_i(0)-v_{\max}\min(t,t_i^c)
        -d_i^e(1-e^{-k[t-t_i^c]_+}).
\end{equation}
The predicted position follows the initial error direction. The corresponding
surge-energy increment has a closed form:
\begin{align}
 \widehat E_i(H)=\frac{18}{3600}\bigg[
 &v_{\max}^3\min(H,t_i^c)\nonumber\\[-1mm]
 &+\frac{(kd_i^e)^3}{3k}
 \left(1-e^{-3k[H-t_i^c]_+}\right)\bigg].\label{eq:energy}
\end{align}
This is a guidance-cost surrogate. The simulator separately executes yaw and
surge lags, acceleration limits, currents, sway disturbances and local repulsion.
Equation (\ref{eq:energy}) excludes common hotel power and does not bound
actual energy in moving water.

\subsection{Common objective and proposal budget}
At one-minute decisions, evaluate each candidate at three forecast times
$H/6,H/2,5H/6$, for $H\leq300$ s:
\begin{equation}
 J(g)=\frac13\sum_{\ell=1}^{3}
 \widehat F(\widehat p(t_\ell);w(t+t_\ell))
 -\lambda\frac{\sum_i\widehat E_i(H)}
 {N(100~\mathrm W)H/3600}.\label{eq:J}
\end{equation}
Only the explicitly scheduled monitoring priority is forecast; future weather
is not supplied. Defaults are $\lambda=0.025$, maximum local displacement
120 m and acceptance margin $\epsilon=0.0005$.
The reference action continues the previously issued waypoints. A new goal is
accepted only if $J(g)>J(g_{\mathrm{previous}})+\epsilon$.
A hold-at-current-cached-position option is also evaluated.

Every new policy evaluates ten candidate waypoint sets with three temporal
scores each. Predicted pair separation must exceed 220 m at the quadrature
positions and horizon endpoint; candidate radii cannot exceed $1.4R$.
This sampled geometric check is not a continuous collision-avoidance proof.
If any coordinator pose is more than 180 s old, previously issued goals persist;
the onboard controller independently holds when its own command expires.

The conventional economic proposals use centered finite differences of the
joint temporal score in each position coordinate, with perturbation 80 m.
They include the normalized gradient, shorter steps, a two-vessel subset,
small rotated directions, and radial expansion/contraction. This is a local,
restricted optimizer, not a globally solved nonlinear MPC problem.

\section{Biological proposal mechanisms}
\subsection{Branching and hydraulic fusion}
The branching direction $g^b$ is a finite-difference gradient of the temporal
sensing score with the radio network made fully connected. It proposes where
sensing could improve before accounting for delivery.

For the hydraulic fusion analogy, use nominal $q_{ij}$ as conductances, ground
the base node, and form the grounded Laplacian $L_g$. Let $b_i$ be normalized
unique instantaneous sensing contribution, with a small positive floor:
\begin{equation}
 b_i\propto 10^{-5}+\sum_m w_m\,a_ip_{im}
                       \prod_{j\ne i}(1-a_jp_{jm}).
\end{equation}
Holding demand $b$ fixed, solve $L_gv=b$ and set $U=b^\top L_g^{-1}b$.
For edge conductance $q_{ij}$,
\begin{equation}
 \frac{\partial U}{\partial q_{ij}}=-(v_i-v_j)^2.\label{eq:hydraulic}
\end{equation}
The base potential is zero. Equation (\ref{eq:hydraulic}) follows by
differentiating the matrix inverse; its spatial gradient follows through
(\ref{eq:q}). Because $b$ itself changes with geometry, the implemented
derivative is explicitly a fixed-demand partial derivative.

The fungal proposal family combines normalized branching and fusion directions:
\begin{equation}
 \delta p^{(f)}=h\,\mathcal N\big[(1-f)\mathcal N(g^b)
                       +f\mathcal N(-\nabla_p U)\big],
\end{equation}
for $f\in\{0,0.2,0.4,0.6,0.8,1\}$, $h=120$ m, and normalization
$\mathcal N(x)=x/\max_i\|x_i\|$ when the denominator is nonzero.
All candidates pass through (\ref{eq:J}). Radio links are not created by a
virtual fungal trail: only physical position changes modify their probabilities.

\subsection{Deadline-dependent fusion}
Hydraulic conductance does not encode report deadlines. Consider a nominal
link probability $q_e$ shared over the planning window, with realized round-$k$
success probability $f_kq_e$. Conditional on shared fade, the link-round
Bernoulli trials are independent. Let $F^{e,k,+}$ and $F^{e,k,-}$ be
(\ref{eq:temporal}) with only that edge-round forced present or absent.

\begin{proposition}[Conditional link influence]
For fixed sensing footprints and visibility law,
\begin{equation}
 I_e:=\frac{\partial \E[F]}{\partial q_e}
 =\sum_{k=1}^{L}\E\left[f_k
                 (F^{e,k,+}-F^{e,k,-})\right]\geq0.\label{eq:influence}
\end{equation}
\end{proposition}
\begin{proof}
Condition on all latent variables and all link-round outcomes except the
selected trial. The conditional expectation is
$f_kq_eF^{e,k,+}+(1-f_kq_e)F^{e,k,-}$.
Differentiate and then average. Since $q_e$ appears in every round, sum the
partial derivatives. Adding an edge cannot remove a feasible causal report
path, hence each difference is nonnegative.
\end{proof}

This is an application of the standard multilinear Bernoulli expectation
identity. It supplies a useful engineering replacement for hydraulic
importance, not a new general reliability theorem. Shared fades remain inside
the conditioning; unconditional independent-link assumptions are unnecessary.

The deadline-fusion direction is
\begin{equation}
 g^{d}_i=\sum_{e\ni i}I_e\,\nabla_{p_i}q_e.
\end{equation}
It replaces $-\nabla_p U$ in the same six mixture proposals. Monte Carlo
interventions estimate (\ref{eq:influence}) using the common planning paths.
This is a \emph{radio-only} spatial derivative: sensing and spatial visibility
correlation also change with motion. Consequently, mixing $g^d$ and $g^b$ is
still a proposal heuristic; the common objective evaluates the full candidate.

\subsection{Recruitment and its conventional control}
The ant-inspired variant starts from the conventional joint-score gradient.
Its normalized per-vessel magnitude is $z_i$. Nominal radio weights, attenuated
by received telemetry age, form a row-normalized matrix $Q$ with zero diagonal.
Activity and refractory state evolve as
\begin{align}
 a_i^+={}&\eta a_i+(1-\eta)
  \left[\tanh(2z_i+\gamma[Qa]_i-0.6r_i)\right]_+,\\
 r_i^+={}&\zeta r_i+(1-\zeta)
                 \1\{\text{waypoint changed for vessel }i\},
\end{align}
where $\eta=e^{-60/120}$, $\zeta=e^{-60/480}$ and $\gamma=0.5$.
The two largest activities receive new excursion proposals; other waypoints
persist. Universal hold and continuation candidates remain available, so this
is not a hard bound of two moving vessels at every physical instant.

The matched conventional sparse policy selects the two largest instantaneous
gradient magnitudes, without activity, coupling or refractory memory. A
no-social comparator retains the memory and refractory equations but sets
$\gamma=0$. This distinguishes a biological interpretation from evidence that
the coupling itself helps.

\subsection{Computational complexity}
One temporal score costs
$O(SL(N^2+NM))$ for $M$ monitoring points.
The local policies use $4N$ finite-difference field calls and 30 trajectory
scores per decision, plus a small $O(N^3)$ grounded solve for hydraulic fusion.
Deadline fusion additionally requires $2L|\mathcal E|$ conditional score
evaluations, giving a direct implementation cost
$O(S|\mathcal E|L^2(N^2+NM))$. With a complete candidate radio graph,
$|\mathcal E|=N(N+1)/2$. This implementation is practical for the tested small
fleet, but is not a demonstrated scalable solver for hundreds of vessels.
Incremental reachability updates, influence screening and sparse graph
representations are concrete routes to reducing that cost.
